\section{Introduction}

Digital interventions and instrumented cognitive tasks can record behavior at a temporal resolution that is lost when outcomes are reduced to session averages or pre--post differences. In many regulatory problems, however, the clinically or mechanistically relevant question is not only how strongly a person reacts to a challenge but how the system evolves afterward.

Emotion-dynamics research already emphasizes temporal characteristics such as variability, inertia, and change processes \parencite{reitsema2022emotiondynamics}. Yet the term \emph{recovery} is used inconsistently across behavioral, affective, psychophysiological, and digital-intervention literatures. A recovery endpoint may refer to latency until a threshold is crossed, slope toward baseline, persistence within a target range, area under a deviation curve, or merely a post-event difference. These definitions are not interchangeable.

In controlled computerized tasks, the operational meaning of post-interruption recovery depends on the observed outcome. A ten-response timecourse in a complex game distinguishes the latency of the first resumed response from subsequent performance stabilization \parencite{altmann2007timecourse}. In a different sequential-choice paradigm, brief interruptions primarily increased first-trial step-order errors, while response latencies were less sensitive or unchanged \parencite{altmann2014momentary}. These illustrate two distinct measurement families rather than interchangeable estimates of a single elapsed time to recovery. Both studies involved university students rather than clinical pediatric populations; their relevance here is conceptual and methodological, pending formal eligibility adjudication.

The present review therefore focuses on measurement architecture rather than intervention efficacy. Its primary question is: how do interactive digital studies operationalize recovery after a discrete perturbation, and what information is required to reproduce that endpoint? A secondary objective is to identify whether recovery is defined relative to an individual reference, population norm, fixed threshold, model-predicted state, or task-specific criterion.

The review follows PRISMA 2020 reporting principles \parencite{page2021prisma} while using a scoping synthesis because the expected evidence base spans heterogeneous tasks, populations, and outcome modalities.
